\documentclass[10pt,a4paper]{amsart}
\usepackage[margin=19mm]{geometry}
\usepackage{amsmath,amssymb,mathtools}
\usepackage[scr=rsfs,cal=euler]{mathalfa}
\usepackage{xfrac}
\usepackage[unicode,hidelinks,bookmarks=false]{hyperref}
\newcommand{\ie}{i.e.\ }
\newcommand{\cf}{cf.\ }
\newcommand{\resp}{respectively\ }
\newcommand{\Subsection}{\subsection}
\providecommand{\nobreakdash}{\nobreak\hbox{-}}
\setlength{\emergencystretch}{2em}
\multlinegap0pt
\usepackage{amssymb}
\usepackage{xcolor}
\usepackage{enumerate}
\usepackage{mathtools}
\newtheorem{lemma}{Lemma}
\title{Instantiating Bayesian CVaR lower bounds in Interactive Decision Making Problems}
\author{Raghav Bongole, Tobias J. Oechtering,, Mikael Skoglund}
\date{Source: arXiv:2604.12519v1 (CC BY 4.0)}
\begin{document}
\maketitle

\noindent\emph{Source and licence.} Instantiating Bayesian CVaR lower bounds in Interactive Decision Making Problems. Version arXiv:2604.12519v1 (CC BY 4.0), distributed by arXiv under the Creative Commons Attribution 4.0 International licence, http://creativecommons.org/licenses/by/4.0/, as recorded in the arXiv licence metadata for this exact version and shown on the abstract page https://arxiv.org/abs/2604.12519v1. Full licence notice: \texttt{licenses/arxiv-2604.12519-CC-BY-4.0.txt}.\par\smallskip

\noindent\textit{Editorial scope.} Counted unit: the article's lemma lem:scalar-min (source0.tex lines 1399--1415). The article's complete original proof of it is delivered with it (source0.tex lines 1417--1585, one \texttt{proof} environment of 169 lines). No mathematics is written, repaired or re-derived: the statement and the proof are the article's own lines, in the article's own notation, and no retained line is substituted or re-flowed.  No further source-local context is needed: every cross-reference of every retained line resolves inside the two delivered spans.  Nothing was omitted from any retained span, so the package carries no omission record.  No retained line contains a citation, so the wrapper carries no bibliography and no bibliography entry.  No retained line contains a figure environment, an image-inclusion command or drawing code, so no image is delivered and the package declares no asset dependency.  Editorial formatting only, in addition: an AMS \texttt{amsart} preamble that restates the article's own declarations and macros used by the retained lines, namely the environments \texttt{lemma} on separate counters, together with the standard packages those lines need.  The retained environments print in this document as Lemma 1 in order of appearance, whereas the article prints the same items with the numbering of its own full document.  The complete native source of the article as distributed in the arXiv e-print for this exact version is bundled unchanged as \texttt{sources/198407/source0.tex}.
\begin{lemma}[Scalar minimization]
\label{lem:scalar-min}
For \(\alpha\in[0,1)\) and \(\rho\ge 0\), define
\[
F_{\alpha,\rho}(x):=
\frac12-x+\frac{(\sqrt{x}-\rho/\sqrt2)_+^2}{1-\alpha},
\qquad x\in[0,1/2].
\]
Then
\[
\inf_{x\in[0,1/2]}F_{\alpha,\rho}(x)=\Psi_\alpha(\rho).
\]
Moreover,
\[
\sup_{\rho\ge 0}\rho\,\Psi_\alpha(\rho)=c_\alpha.
\]
\end{lemma}

\begin{proof}
Write \(x=s^2\), with \(s\in[0,1/\sqrt2]\). Then
\[
F_{\alpha,\rho}(x)
=
\frac12-s^2+\frac{(s-\rho/\sqrt2)_+^2}{1-\alpha}.
\]

We first identify \(\inf_{x\in[0,1/2]}F_{\alpha,\rho}(x)\).

If \(0\le s\le \rho/\sqrt2\), then
\[
F_{\alpha,\rho}(x)=\frac12-s^2,
\]
so the minimum on this region is
\[
\begin{cases}
\frac12-\frac{\rho^2}{2}, & 0\le \rho\le 1,\\[1ex]
0, & \rho\ge 1.
\end{cases}
\]

Now suppose \(\rho/\sqrt2\le s\le 1/\sqrt2\). Then
\[
F_{\alpha,\rho}(x)
=
\frac12-s^2+\frac{(s-\rho/\sqrt2)^2}{1-\alpha}.
\]

If \(\alpha=0\), this becomes
\[
\frac12-\sqrt2\,\rho\,s+\frac{\rho^2}{2},
\]
which is decreasing in \(s\), hence minimized at \(s=1/\sqrt2\), with value
\[
\frac{(1-\rho)^2}{2}
\qquad (0\le \rho\le 1).
\]
Since
\[
\frac{(1-\rho)^2}{2}\le \frac12-\frac{\rho^2}{2}
\qquad (0\le \rho\le 1),
\]
it follows that
\[
\inf_{x\in[0,1/2]}F_{0,\rho}(x)
=
\begin{cases}
\frac{(1-\rho)^2}{2}, & 0\le \rho\le 1,\\[1ex]
0, & \rho\ge 1,
\end{cases}
=
\Psi_0(\rho).
\]

Assume now \(0<\alpha<1\). On \([\rho/\sqrt2,1/\sqrt2]\),
\[
G_{\alpha,\rho}(s):=
\frac12-s^2+\frac{(s-\rho/\sqrt2)^2}{1-\alpha}
\]
is a convex quadratic with
\[
G'_{\alpha,\rho}(s)=\frac{2\alpha}{1-\alpha}s-\frac{\sqrt2\,\rho}{1-\alpha},
\]
so its critical point is
\[
s^\star=\frac{\rho}{\sqrt2\,\alpha}.
\]

If \(0\le \rho\le \alpha\), then \(s^\star\in[\rho/\sqrt2,1/\sqrt2]\), and
\[
G_{\alpha,\rho}(s^\star)=\frac12-\frac{\rho^2}{2\alpha}.
\]
If \(\alpha<\rho\le 1\), then \(s^\star>1/\sqrt2\), so \(G_{\alpha,\rho}\) is decreasing on \([\rho/\sqrt2,1/\sqrt2]\), and its minimum is attained at \(s=1/\sqrt2\), giving
\[
G_{\alpha,\rho}(1/\sqrt2)=\frac{(1-\rho)^2}{2(1-\alpha)}.
\]
Finally, if \(\rho\ge 1\), choosing \(s=1/\sqrt2\) gives value \(0\).

Comparing with the first region,
\[
\frac12-\frac{\rho^2}{2\alpha}\le \frac12-\frac{\rho^2}{2}
\qquad (0\le \rho\le \alpha),
\]
and
\[
\frac{(1-\rho)^2}{2(1-\alpha)}\le \frac12-\frac{\rho^2}{2}
\qquad (\alpha<\rho\le 1).
\]
Hence
\[
\inf_{x\in[0,1/2]}F_{\alpha,\rho}(x)=\Psi_\alpha(\rho)
\]
for all \(\alpha\in[0,1)\).

It remains to maximize \(\rho\Psi_\alpha(\rho)\). Since \(\Psi_\alpha(\rho)=0\) for \(\rho\ge 1\), it suffices to work on \([0,1]\).

If \(\alpha=0\), then
\[
\rho\Psi_0(\rho)=\frac{\rho(1-\rho)^2}{2},\qquad 0\le \rho\le 1,
\]
whose derivative is
\[
\frac{(1-\rho)(1-3\rho)}{2}.
\]
Thus the maximum is attained at \(\rho=1/3\), with value
\[
\frac{(1/3)(2/3)^2}{2}=\frac{2}{27}=c_0.
\]

Now let \(0<\alpha<1\). We have
\[
\rho\Psi_\alpha(\rho)=
\begin{cases}
\displaystyle
\frac{\rho}{2}-\frac{\rho^3}{2\alpha}, & 0\le \rho\le \alpha,\\[2ex]
\displaystyle
\frac{\rho(1-\rho)^2}{2(1-\alpha)}, & \alpha<\rho\le 1.
\end{cases}
\]

On \([0,\alpha]\), the derivative is
\[
\frac12-\frac{3\rho^2}{2\alpha},
\]
so the critical point is \(\rho=\sqrt{\alpha/3}\). Hence the maximum on \([0,\alpha]\) is
\[
\frac{\sqrt{\alpha}}{3\sqrt3}
\quad\text{if }\alpha\ge \frac13,
\]
and otherwise is attained at \(\rho=\alpha\), with value
\[
\frac{\alpha(1-\alpha)}{2}.
\]

On \((\alpha,1]\), the derivative is
\[
\frac{(1-\rho)(1-3\rho)}{2(1-\alpha)},
\]
so the interior critical point is \(\rho=1/3\). Hence the maximum on \((\alpha,1]\) is
\[
\frac{2}{27(1-\alpha)}
\quad\text{if }\alpha\le \frac13,
\]
and otherwise is attained at the boundary \(\rho=\alpha\), with value
\[
\frac{\alpha(1-\alpha)}{2}.
\]

Since the two branches agree at \(\rho=\alpha\), it remains only to compare the boundary value with the interior one. For \(0\le \alpha\le 1/3\),
\[
\frac{\alpha(1-\alpha)}{2}\le \frac{2}{27(1-\alpha)},
\]
and for \(1/3\le \alpha<1\),
\[
\frac{\alpha(1-\alpha)}{2}\le \frac{\sqrt{\alpha}}{3\sqrt3}.
\]
Therefore
\[
\sup_{\rho\ge 0}\rho\,\Psi_\alpha(\rho)=
\begin{cases}
\displaystyle
\frac{2}{27(1-\alpha)}, & 0\le \alpha\le \frac13,\\[2ex]
\displaystyle
\frac{\sqrt{\alpha}}{3\sqrt3}, & \frac13\le \alpha<1,
\end{cases}
\]
which is exactly \(c_\alpha\).
\end{proof}

\end{document}
